% Imagen y núcleo

\begin{frame}
    \frametitle{Núcleo e imagen}

    \begin{definition}
        Sea $f\colon V\to V'$ un homomorfismo. El \alert{núcleo} y la
        \alert{imagen} de $f$ son
        \[
            \ker f=\{u\in V:\ f(u)=0\}\subseteq V,
            \qquad
            \im f=\{f(u):\ u\in V\}\subseteq V'.
        \]
    \end{definition}

    \medskip

    El núcleo recoge lo que $f$ \alert{anula}, y la imagen lo que $f$
    \alert{alcanza}.

    \medskip

    La aplicación $f$ es sobreyectiva si y solo si $\im f=V'$.
\end{frame}


\begin{frame}
    \frametitle{El núcleo y la imagen son subespacios}

    \begin{proposition}
        Sea $f\colon V\to V'$ un homomorfismo. Entonces $\ker f$ es un
        subespacio de $V$ e $\im f$ es un subespacio de $V'$.
    \end{proposition}

    \begin{proof}
        \alert{Núcleo.} $f(0)=0$, luego $0\in\ker f$. Si $u,u'\in\ker f$ y
        $\lambda,\mu\in k$,
        \[
            f(\lambda u+\mu u')=\lambda f(u)+\mu f(u')=\lambda 0+\mu 0=0,
        \]
        luego $\lambda u+\mu u'\in\ker f$.

        \medskip

        \alert{Imagen.} $0=f(0)\in\im f$. Si $w=f(u)$ y $w'=f(u')$,
        entonces
        \[
            \lambda w+\mu w'=\lambda f(u)+\mu f(u')=f(\lambda u+\mu u')\in\im f.\qedhere
        \]
    \end{proof}
\end{frame}


\begin{frame}
    \frametitle{Ejemplos de núcleo e imagen}

    \begin{example}
        \begin{enumerate}
            \item $f_A\colon k^n\to k^m$, $u\mapsto Au$:
            $\ker f_A=\{u:\ Au=0\}$ es el conjunto de soluciones del
            sistema homogéneo, e $\im f_A=\{Au:\ u\in k^n\}$ es el
            subespacio generado por las columnas de $A$.
            \item La proyección $\R^3\rightarrow\R^3$, $(u_1,u_2,u_3)\mapsto(u_1,u_2,0)$,
            tiene núcleo $\langle e_3\rangle$ e imagen el plano
            $u_3=0$.
            \item La aplicación nula $0\colon V\to V'$ tiene $\ker=V$ e $\im=\{0\}$. La
            identidad $\id\colon V\to V$ tiene $\ker=\{0\}$ e $\im=V$.
            \item La rotación $f_\theta$ de $\R^2$ tiene $\ker f_\theta=\{0\}$
            e $\im f_\theta=\R^2$.
            \item La proyección canónica $\pi\colon V\to V/U$ tiene
            $\ker\pi=U$ (pues $\pi(u)=u+V=0+V$ si y solo si
            $u\in U$) e $\im\pi=V/U$.
        \end{enumerate}
    \end{example}
\end{frame}


\begin{frame}
    \frametitle{Más ejemplos de núcleo e imagen}

    \begin{example}
        \begin{enumerate}
            \item $D\colon\R[x]\to\R[x]$, $D(p)=p'$, $\ker D$ son los
            polinomios constantes e $\im D=\R[x]$ (todo polinomio tiene
            una primitiva). Si restringimos a $D\colon\R_n[x]\to\R_{n-1}[x]$,
            sigue siendo $\ker D=\R$ e $\im D=\R_{n-1}[x]$.
            \item $\operatorname{ev}_a\colon k[x]\to k$, $p\mapsto p(a)$,
            $\ker\operatorname{ev}_a=\{p(x)\in k[x]:\ p(a)=0\}$, los polinomios que tienen a $a$ como raíz, e
            $\im\operatorname{ev}_a=k$ (la constante $\lambda$ es
            $\operatorname{ev}_a(\lambda)$).
            \item La multiplicación por $x$, $k[x]\to k[x]$, $p(x)\mapsto xp(x)$, tiene núcleo $\{0\}$
            e imagen $\{p(x)\in k[x]:\ p(0)=0\}$, los polinomios con término independiente nulo.
            \item $\operatorname{tr}\colon M_n(k)\to k$,
            $\ker\operatorname{tr}$ son las matrices de traza nula e
            $\im\operatorname{tr}=k$.
            \item La trasposición $M_{m\times n}(k)\to M_{n\times m}(k)$, $A\mapsto A^t$, tiene núcleo $\{0\}$ e
            imagen $M_{n\times m}(k)$.
            \item La asignación de coordenadas $c_B\colon V\to k^n$, núcleo $\{0\}$ e imagen $k^n$.
        \end{enumerate}
    \end{example}
\end{frame}


\begin{frame}
    \frametitle{Inyectividad y núcleo}

    \begin{proposition}
        Sea $f\colon V\to V'$ un homomorfismo. Para todo $u,u'\in V$,
        \[
            f(u)=f(u')\iff u-u'\in\ker f.
        \]
        En particular, $f$ es inyectiva si y solo si $\ker f=\{0\}$.
    \end{proposition}

    \begin{proof}
        $f(u)=f(u')\iff f(u-u')=f(u)-f(u')=0\iff u-u'\in\ker f$.

        Si $f$ es inyectiva y $u\in \ker f$ entonces $f(u)=0=f(0)$, luego $u=0$, por tanto $\ker f=\{0\}$.
        
        Recíprocamente,
        si $\ker f=\{0\}$ y $f(u)=f(u')$, entonces $u-u'\in\ker f=\{0\}$,
        luego $u-u'=0$, es decir, $u=u'$, así que $f$ es inyectiva.
    \end{proof}

    \medskip

    Para estudiar la inyectividad de un homomorfismo basta, por tanto,
    resolver una \alert{única} ecuación, $f(u)=0$.
\end{frame}


\begin{frame}
    \frametitle{La imagen está generada por las imágenes de los generadores}

    \begin{proposition}
        Sea $f\colon V\to V'$ un homomorfismo y sean $u_1,\ldots,u_r$
        generadores de $V$. Entonces
        \[
            \im f=\langle f(u_1),\ldots,f(u_r)\rangle.
        \]
    \end{proposition}

    \begin{proof}
        Probamos las dos inclusiones.

        ($\supseteq$) Cada $f(u_j)$ está en $\im f$ por tanto 
        $\langle f(u_1),\ldots,f(u_r)\rangle\subseteq\im f$.

        ($\subseteq$) Sea $w\in\im f$, es decir, $w=f(u)$ para algún
        $u\in V$. Como $u_1,\ldots,u_r$ generan $V$, existen
        $\lambda_1,\ldots,\lambda_r\in k$ con
        $u=\lambda_1u_1+\cdots+\lambda_ru_r$. Como $f$ preserva
        combinaciones lineales,
        \[
            w=f(u)=\lambda_1f(u_1)+\cdots+\lambda_rf(u_r)
            \in\langle f(u_1),\ldots,f(u_r)\rangle.\qedhere
        \]
    \end{proof}

    \medskip

    En particular, si $V$ tiene dimensión finita, $\im f$ también.
\end{frame}


\begin{frame}
    \frametitle{Imagen de familias de vectores}

    \begin{proposition}
        Sea $f\colon V\to V'$ un homomorfismo y sean $u_1,\ldots,u_r\in V$.
        \begin{enumerate}[(1)]
            \item Si $f(u_1),\ldots,f(u_r)$ son linealmente independientes,
            también lo son $u_1,\ldots,u_r$.
            \item Si $f$ es inyectiva y $u_1,\ldots,u_r$ son linealmente
            independientes, también lo son $f(u_1),\ldots,f(u_r)$.
            \item Si $f$ es sobreyectiva y $u_1,\ldots,u_r$ generan $V$,
            entonces $f(u_1),\ldots,f(u_r)$ generan $V'$.
            \item Si $f$ es un isomorfismo y $B$ es una base de $V$,
            entonces $f(B)$ es una base de $V'$.
        \end{enumerate}
    \end{proposition}

    \medskip

    Un homomorfismo cualquiera \alert{no} conserva la independencia
    lineal: la aplicación nula lleva cualquier familia no vacía a una familia de
    vectores nulos, que es linealmente dependiente.
\end{frame}


\begin{frame}
    \frametitle{Imagen de familias de vectores: demostración}

    \begin{proof}
        \alert{(1)} Si $\sum_{i=1}^r\lambda_iu_i=0$, aplicando $f$ se obtiene
        $\sum_{i=1}^r\lambda_if(u_i)=0$, luego todos los $\lambda_i$ son nulos
        por la independencia lineal de $f(u_1),\ldots,f(u_r)$.

        \alert{(2)} Si $\sum_{i=1}^r\lambda_if(u_i)=0$, entonces
        $f\bigl(\sum_{i=1}^r\lambda_iu_i\bigr)=\sum_{i=1}^r\lambda_if(u_i)=0$, y como $\ker f=\{0\}$,
        $\sum_{i=1}^r\lambda_iu_i=0$. Por independencia lineal, todos los $\lambda_i$
        son nulos.

        \alert{(3)} Por la proposición anterior,
        $\langle f(u_1),\ldots,f(u_r)\rangle=\im f=V'$.

        \alert{(4)} Un isomorfismo es inyectivo y sobreyectivo, así que
        (2) y (3) dan que $f(B)$ es independiente y genera $V'$.
    \end{proof}
\end{frame}


\begin{frame}
    \frametitle{Un ejemplo antes del teorema}

    \begin{columns}[c]
        \begin{column}{0.55\textwidth}
            \begin{example}
                Sea $f\colon\R^2\to\R$, $f(u_1,u_2)=u_1-u_2$. Es sobreyectiva
                y $\ker f=\langle(1,1)\rangle$.
            \end{example}

            Para cada $c\in\R$, la \alert{fibra}
            \[
                f^{-1}(c)=\{u:\ u_1-u_2=c\}
            \]
            es una recta paralela a $\ker f$, es decir, una clase módulo
            $\ker f$.

            \medskip

            Cada número real $c$ corresponde a una única clase de
            $\R^2/\ker f$, y viceversa: $f$ identifica el cociente con
            $\R=\im f$.
        \end{column}
        \begin{column}{0.42\textwidth}
            \centering
            \begin{tikzpicture}[scale=0.7]
                \clip (-3.2,-2.6) rectangle (3.4,2.8);
                \draw[->] (-2.5,0) -- (2.5,0) node[below left] {\footnotesize $u_1$};
                \draw[->] (0,-2.5) -- (0,2.5) node[below left] {\footnotesize $u_2$};
                \foreach \c in {-2,-1,1,2}
                    \draw[gray!60] (-3,-3-\c) -- (3,3-\c);
                \draw[very thick,alertcolor] (-3,-3) -- (3,3);
                \foreach \c in {0,1,2}
                    \node[right,font=\footnotesize,fill=white,inner sep=1pt] at (2.5,2.5-\c) {$\c$};
                \foreach \c in {-2,-1}
                    \node[left,font=\footnotesize,fill=white,inner sep=1pt] at (-2.5,-2.5-\c) {$\c$};
            \end{tikzpicture}
        \end{column}
    \end{columns}
\end{frame}


\begin{frame}
    \frametitle{Las fibras son las clases módulo el núcleo}

    \begin{lemma}
        Sea $f\colon V\to V'$ un homomorfismo. Para todo $u,u'\in V$,
        \[
            f(u)=f(u')\iff u+\ker f=u'+\ker f.
        \]
    \end{lemma}

    \begin{proof}
        Por la caracterización de la inyectividad,
        $f(u)=f(u')\iff u-u'\in\ker f$, y por la caracterización de las
        clases módulo un subespacio (Tema~3),
        $u-u'\in\ker f\iff u+\ker f=u'+\ker f$.
    \end{proof}

    \medskip

    Es decir, $f$ \alert{no distingue} los vectores de una misma clase
    módulo $\ker f$, y distingue los de clases distintas.
\end{frame}


\begin{frame}[allowframebreaks, t]
    \frametitle{Primer teorema de isomorfía}

    \begin{theorem}
        Sea $f\colon V\to V'$ un homomorfismo. La aplicación
        \[
            \bar f\colon V/\ker f\longrightarrow\im f,
            \qquad
            \bar f(u+\ker f)=f(u),
        \]
        está bien definida y es un isomorfismo. En particular,
        \[
            V/\ker f\cong\im f.
        \]
    \end{theorem}

    \framebreak

    \begin{proof}
        Escribamos $K=\ker f$. Por el lema, para todo $u,u'\in V$,
        \[
            u+K=u'+K\iff f(u)=f(u').
        \]
        La implicación ($\Rightarrow$) dice que $\bar f$ \alert{está bien
        definida}: el valor $f(u)$ no depende del representante $u$ de la
        clase. La implicación ($\Leftarrow$) dice que $\bar f$ es
        \alert{inyectiva}.

        \alert{Lineal.} Para $u,u'\in V$ y $\lambda\in k$,
        \[
            \bar f\bigl((u+K)+(u'+K)\bigr)=\bar f\bigl((u+u')+K\bigr)
            =f(u+u')=\bar f(u+K)+\bar f(u'+K),
        \]
        \[
            \bar f\bigl(\lambda(u+K)\bigr)=\bar f(\lambda u+K)
            =f(\lambda u)=\lambda\bar f(u+K).
        \]

        \alert{Sobreyectiva.} Todo elemento de $\im f$ es $f(u)=\bar f(u+K)$
        para algún $u\in V$.
    \end{proof}
\end{frame}


\begin{frame}[fragile]
    \frametitle{Factorización canónica de un homomorfismo}

    Todo homomorfismo $f\colon V\to V'$ se descompone como
    \[
        \begin{tikzcd}[column sep=large]
            V\arrow[rr,"f"]\arrow[dr,"\pi"'] & & V'\\
            & V/\ker f\arrow[r,"\bar f"',"\cong"] & \im f\arrow[u,hook,"\iota"']
        \end{tikzcd}
    \]
    es decir, $f=\iota\circ\bar f\circ\pi$, donde
    \begin{itemize}
        \item $\pi$ es la proyección canónica, \alert{sobreyectiva};
        \item $\bar f$ es un \alert{isomorfismo};
        \item $\iota$ es la inclusión, \alert{inyectiva}.
    \end{itemize}

    \medskip

    En efecto, $\iota\bar f\pi(u)=\iota\bar f(u+\ker f)=f(u)$.

    \medskip

    Así, \alert{todo} homomorfismo es, salvo isomorfismo, una proyección
    sobre un cociente seguida de una inclusión.
\end{frame}


\begin{frame}
    \frametitle{Teorema de la dimensión}

    \begin{corollary}
        Sea $f\colon V\to V'$ un homomorfismo con $V$ de dimensión
        finita. Entonces
        \[
            \dim V=\dim\ker f+\dim\im f.
        \]
    \end{corollary}

    \begin{proof}
        Por el primer teorema de isomorfía, $V/\ker f\cong\im f$, y los
        espacios isomorfos de dimensión finita tienen la misma dimensión.
        Por tanto,
        \[
            \dim\im f=\dim V/\ker f=\dim V-\dim\ker f
        \]
        (Tema~3), es decir, $\dim V=\dim\ker f+\dim\im f$.
    \end{proof}

    \medskip

    El teorema es una versión abstracta de
    $\dim\{Ax=0\}=n-\rg A$ (Tema~3).
\end{frame}


\begin{frame}
    \frametitle{Núcleo e imagen en coordenadas}

    \begin{proposition}
        Sea $f\colon V\to V'$ un homomorfismo, $B$ y $B'$ bases de $V$ y
        $V'$, y $A=M_{B,B'}(f)$. Entonces
        \begin{enumerate}[(1)]
            \item $u\in\ker f\iff Au_B=0$;
            \item $w\in\im f\iff w_{B'}$ es combinación lineal de las
            columnas de $A$.
        \end{enumerate}
        Por tanto, $\dim\im f=\rg A$.
    \end{proposition}

    \begin{proof}
        \alert{(1)} $f(u)=0\iff f(u)_{B'}=0\iff Au_B=0$, porque $c_{B'}$ es
        inyectiva y $0_{B'}=0$.

        \alert{(2)} Si $w=f(u)$, entonces $w_{B'}=Au_B$ es una combinación
        lineal de las columnas de $A$. Recíprocamente, si $w_{B'}=Ax$,
        como $c_B$ es sobreyectiva, $x=u_B$ para algún $u\in V$, y
        $f(u)_{B'}=Ax=w_{B'}$, luego $w=f(u)\in\im f$.

        Por (2), $c_{B'}$ lleva $\im f$ isomórficamente sobre el espacio
        generado por las columnas de $A$, cuya dimensión es $\rg A$
        (Tema~3).
    \end{proof}
\end{frame}


\begin{frame}
    \frametitle{Cálculo de bases del núcleo y de la imagen}

    Sea $A=M_{B,B'}(f)$, con $B=\{v_1,\ldots,v_n\}$ y $B'=\{v'_1,\ldots,v'_m\}$.
    Escalonamos $A$.

    \begin{itemize}
        \item \alert{Núcleo.} Resolvemos $Ax=0$ (Tema~1). Si
        $x^{(1)},\ldots,x^{(s)}$ es una base de las soluciones,
        \[
            u^{(i)}=x^{(i)}_1v_1+\cdots+x^{(i)}_nv_n,
            \qquad i=1,\ldots,s,
        \]
        es una base de $\ker f$, y $s=n-\rg A$.
        \item \alert{Imagen.} Si las columnas $j_1,\ldots,j_r$ de una forma
        escalonada de $A$ tienen pivote, los vectores
        \[
            f(v_{j_1}),\ \ldots,\ f(v_{j_r})
        \]
        (cuyas coordenadas en $B'$ son esas columnas de $A$) forman una
        base de $\im f$, y $r=\rg A$.
    \end{itemize}

    \medskip

    Se trabaja en coordenadas con los métodos de los Temas~1 y~3, y se
    vuelve a $V$ y a $V'$ mediante $c_B^{-1}$ y $c_{B'}^{-1}$.
\end{frame}


\begin{frame}[allowframebreaks, t]
    \frametitle{Ejemplo: núcleo e imagen de un homomorfismo de $\R^4$ en $\R^3$}

    \begin{example}
        Sea $f\colon\R^4\to\R^3$,
        \[
            f(u_1,u_2,u_3,u_4)=(u_1+2u_2+u_4,\ 2u_1+4u_2+u_3+3u_4,\ u_1+2u_2+u_3+2u_4).
        \]
        Calculemos bases de $\ker f$ e $\im f$.
    \end{example}

    Con las bases canónicas, la matriz de $f$ es
    \[
        A=\begin{pmatrix}
            1&2&0&1\\2&4&1&3\\1&2&1&2
        \end{pmatrix}
        \xrightarrow[\;F_3-F_1\;]{\;F_2-2F_1\;}
        \begin{pmatrix}
            1&2&0&1\\0&0&1&1\\0&0&1&1
        \end{pmatrix}
        \xrightarrow{\;F_3-F_2\;}
        \begin{pmatrix}
            \alert{1}&2&0&1\\0&0&\alert{1}&1\\0&0&0&0
        \end{pmatrix}.
    \]

    \framebreak

    Hay pivotes en las columnas $1$ y $3$, luego $\dim\im f=2$.

    \medskip

    \alert{Imagen.} Una base es $\{f(e_1),f(e_3)\}=\{(1,2,1),(0,1,1)\}$,
    las columnas $1$ y $3$ de $A$.

    \medskip

    \alert{Núcleo.} El sistema $Au=0$ equivale a $u_1+2u_2+u_4=0$,
    $u_3+u_4=0$. Con $u_2=s$ y $u_4=t$ libres,
    \[
        u=(-2s-t,\ s,\ -t,\ t)=s(-2,1,0,0)+t(-1,0,-1,1),
    \]
    así que $\{(-2,1,0,0),(-1,0,-1,1)\}$ es una base de $\ker f$.

    \medskip

    Comprobación del teorema de la dimensión: $4=2+2$.
\end{frame}


\begin{frame}
    \frametitle{Ejemplo: la diferencia finita de polinomios}

    \begin{example}
        Sea $\Delta\colon\R_2[x]\to\R_2[x]$, $\Delta(p)(x)=p(x+1)-p(x)$.
        Calculemos $\ker\Delta$ e $\im\Delta$.
    \end{example}

    Como $\Delta(1)=0$, $\Delta(x)=1$ y $\Delta(x^2)=1+2x$, la matriz
    en $B=\{1,x,x^2\}$ es
    \[
        A=\begin{pmatrix}0&1&1\\0&0&2\\0&0&0\end{pmatrix},
    \]
    que ya está escalonada, con pivotes en las columnas $2$ y $3$.

    \medskip

    \alert{Núcleo.} $Au=0$ da $u_2=u_3=0$ con $u_1$ libre:
    $\ker\Delta=\langle1\rangle$, los polinomios constantes.

    \medskip

    \alert{Imagen.} Una base es $\{\Delta(x),\Delta(x^2)\}=\{1,\,1+2x\}$,
    luego $\im\Delta=\langle1,x\rangle=\R_1[x]$.

    \medskip

    Se cumple $3=1+2$.
\end{frame}


\begin{frame}
    \frametitle{Ejemplo: el conmutador con una matriz diagonal}

    \begin{example}
        Sea $T=\begin{pmatrix}1&0\\0&2\end{pmatrix}$ y
        $f\colon M_2(\R)\to M_2(\R)$, $f(X)=TX-XT$. Calculemos su núcleo
        y su imagen.
    \end{example}

    Si $X=\begin{pmatrix}a&b\\c&d\end{pmatrix}$, entonces
    \[
        TX-XT=\begin{pmatrix}a&b\\2c&2d\end{pmatrix}
        -\begin{pmatrix}a&2b\\c&2d\end{pmatrix}
        =\begin{pmatrix}0&-b\\c&0\end{pmatrix}.
    \]
    Por tanto, $f(X)=0$ si y solo si $b=c=0$, y
    \[
        \ker f=\{\text{matrices diagonales}\}=\langle E_{11},E_{22}\rangle,
        \qquad
        \im f=\langle E_{12},E_{21}\rangle.
    \]

    \medskip

    Aquí $\dim\ker f=\dim\im f=2$ y $M_2(\R)=\ker f\oplus\im f$. Pero esto
    \alert{no} ocurre en general: para $D\colon\R_2[x]\to\R_2[x]$,
    $\ker D=\langle1\rangle\subseteq\im D=\langle1,x\rangle$.
\end{frame}


\begin{frame}
    \frametitle{Inyectividad, sobreyectividad y dimensión}

    \begin{corollary}
        Sea $f\colon V\to V'$ un homomorfismo entre espacios de dimensión
        finita.
        \begin{enumerate}[(1)]
            \item Si $f$ es inyectiva, $\dim V\leq\dim V'$. Si $f$ es
            sobreyectiva, $\dim V\geq\dim V'$.
            \item Si $\dim V=\dim V'$, son equivalentes: $f$ es
            inyectiva; $f$ es sobreyectiva; $f$ es un isomorfismo.
        \end{enumerate}
    \end{corollary}

    \begin{proof}
        \alert{(1)} Si $f$ es inyectiva, $\ker f=\{0\}$ y
        $\dim V=\dim\im f\leq\dim V'$. Si es sobreyectiva,
        $\dim V=\dim\ker f+\dim V'\geq\dim V'$.

        \alert{(2)} Sea $n=\dim V=\dim V'$. Entonces
        \begin{align*}
            f\text{ inyectiva}&\iff\dim\ker f=0\iff\dim\im f=n\\
            &\iff\im f=V'\iff f\text{ sobreyectiva},
        \end{align*}
        donde usamos que un subespacio de $V'$ de dimensión $n=\dim V'$
        es $V'$ (Tema~3).
    \end{proof}
\end{frame}


\begin{frame}
    \frametitle{Consecuencias y contraejemplos en dimensión infinita}

    \begin{example}
        \begin{enumerate}
            \item Ningún homomorfismo $\R^3\to\R^2$ es inyectivo, y ninguno
            $\R^2\to\R^3$ es sobreyectivo.
            \item Para ver que un endomorfismo de $k^n$, de $k_n[x]$ o de
            $M_n(k)$ es un isomorfismo, \alert{basta} ver que es
            inyectivo: por ejemplo, $\Delta\colon\R_2[x]\to\R_2[x]$ no es
            un isomorfismo porque $\Delta(1)=0$.
            \item En dimensión infinita (2) \alert{falla}:
            $D\colon\R[x]\to\R[x]$ es sobreyectiva pero no inyectiva, y
            $p\mapsto xp$ es inyectiva pero no sobreyectiva.
        \end{enumerate}
    \end{example}
\end{frame}


\begin{frame}
    \frametitle{Clasificación de los espacios de dimensión finita}

    \begin{theorem}
        Sean $V$ y $V'$ espacios vectoriales de dimensión finita sobre
        $k$. Entonces
        \[
            V\cong V'\iff\dim V=\dim V'.
        \]
        En particular, si $\dim V=n$, entonces $V\cong k^n$.
    \end{theorem}

    \begin{proof}
        ($\Rightarrow$) Un isomorfismo lleva una base de $V$ a una base de
        $V'$ con el mismo número de vectores.

        ($\Leftarrow$) Sean $B$ y $B'$ bases de $V$ y $V'$, ambas con
        $n$ vectores. Las aplicaciones de coordenadas son isomorfismos
        $c_B\colon V\to k^n$ y $c_{B'}\colon V'\to k^n$, luego
        $c_{B'}^{-1}\circ c_B\colon V\to V'$ es un isomorfismo.
    \end{proof}

    \medskip

    La dimensión es el \alert{único} invariante de un espacio vectorial de
    dimensión finita salvo isomorfismo: por ejemplo,
    $\R_3[x]\cong M_2(\R)\cong\R^4$.
\end{frame}


\begin{frame}
    \frametitle{Las fibras de un homomorfismo}

    \begin{proposition}
        Sea $f\colon V\to V'$ un homomorfismo y $w\in V'$.
        \begin{itemize}
            \item Si $w\notin\im f$, la ecuación $f(u)=w$ no tiene solución.
            \item Si $w=f(u_0)\in\im f$, el conjunto de soluciones es
            \[
                f^{-1}(w)=\{u\in V:\ f(u)=w\}=u_0+\ker f.
            \]
        \end{itemize}
    \end{proposition}

    \begin{proof}
        Si $f(u_0)=w$, entonces
        \[
            f(u)=w\iff f(u)=f(u_0)\iff u-u_0\in\ker f
            \iff u\in u_0+\ker f.\qedhere
        \]
    \end{proof}

    \medskip

    Es el resultado del Tema~1: \alert{la solución general de un sistema
    es una solución particular más la solución general del sistema
    homogéneo}. Las fibras de $f$ son las clases módulo $\ker f$.
\end{frame}


\begin{frame}
    \frametitle{Ejemplo: resolver $f(u)=w$}

    \begin{example}
        Sea $f\colon\R^4\to\R^3$ el homomorfismo de antes. Resolvamos
        $f(u)=(1,3,2)$.
    \end{example}

    Escalonamos la matriz ampliada $(A\mid w)$:
    \[
        \left(\begin{array}{cccc|c}
            1&2&0&1&1\\2&4&1&3&3\\1&2&1&2&2
        \end{array}\right)
        \xrightarrow[\;F_3-F_1\;]{\;F_2-2F_1\;}
        \left(\begin{array}{cccc|c}
            1&2&0&1&1\\0&0&1&1&1\\0&0&1&1&1
        \end{array}\right)
        \xrightarrow{\;F_3-F_2\;}
        \left(\begin{array}{cccc|c}
            1&2&0&1&1\\0&0&1&1&1\\0&0&0&0&0
        \end{array}\right).
    \]
    Una solución particular es $u_0=(1,0,1,0)=e_1+e_3$ (pues
    $f(e_1)+f(e_3)=(1,2,1)+(0,1,1)=(1,3,2)$), y por tanto
    \[
        f^{-1}(1,3,2)=(1,0,1,0)+\langle(-2,1,0,0),(-1,0,-1,1)\rangle.
    \]

    \medskip

    En cambio, $f(u)=(1,0,0)$ no tiene solución, porque
    $(1,0,0)\notin\im f=\langle(1,2,1),(0,1,1)\rangle$.
\end{frame}


\begin{frame}[allowframebreaks, t]
    \frametitle{Todo subespacio es un núcleo y una imagen}

    \begin{proposition}
        Sea $U$ un subespacio de $V$, con $V$ de dimensión finita.
        \begin{enumerate}[(1)]
            \item $U$ es el \alert{núcleo} de un homomorfismo: $U=\ker\pi$,
            donde $\pi\colon V\to V/U$ es la proyección canónica.
            \item $U$ es la \alert{imagen} de un homomorfismo inyectivo: si
            $\{u_1,\ldots,u_m\}$ es una base de $U$,
            \[
                g\colon k^m\to V,\qquad
                g(\lambda_1,\ldots,\lambda_m)=\lambda_1u_1+\cdots+\lambda_mu_m,
            \]
            cumple $\im g=U$ y $\ker g=\{0\}$.
        \end{enumerate}
    \end{proposition}

    \begin{proof}
        (1) es el ejemplo de $\pi$. En (2), $\im g=\langle u_1,\ldots,u_m\rangle=U$
        y, como $u_1,\ldots,u_m$ son independientes, $g(\lambda)=0$ implica
        $\lambda=0$.
    \end{proof}

    \medskip

    Así, una \alert{parametrización} de $U$ es un homomorfismo inyectivo
    con imagen $U$ (ecuaciones paramétricas), y unas \alert{ecuaciones
    implícitas} son las que definen $U$ como núcleo.
\end{frame}
